\section{From Axis Mappings to Attribute Programs}
\label{sec:method}

Section~\ref{sec:mathematical-synthesis} defines the complete constructor.
This section compiles that same constructor into a symbolic metadata
program, before a match supplies concrete ranks or axis lists. Axis
mappings locate groups; partitioning separates their roles. The generated
program retains any group computations that are not simplified during
compilation. Native execution and checked symbolic lowering consume
the same program; they do not implement separate complete constructors.

\subsection{Symbolic Axis Mapping}
\label{sec:symbolic-axis-mapping}
Section~\ref{sec:mathematical-synthesis} determines which logical axes each
RHS operation uses. Compilation translates these choices into expressions
for dimension positions, evaluated on each match's metadata to obtain
concrete RHS axis attributes.

A mapping associates a group $\vec G$ with a position-list expression
$p_G$ in tensor $t$:
\begin{equation}
 \gath(\operatorname{Ord}_t,p_G)=\vec G.
 \label{eq:mapping}
\end{equation}
Here $\gath$ selects entries from the logical-axis list
$\operatorname{Ord}_t$ at positions $p_G$. TEPL derives mappings from
the LHS operation declarations and captured axis fields. Intermediate
RHS axis orders follow Section~\ref{sec:mathematical-synthesis}'s construction.

For example, if $\operatorname{Ord}_T=[i_0,i_1,j_0,j_1]$, then
$\vec I=[i_0,i_1]$ has positions $p_I=[0,1]$, and
$\vec J=[j_0,j_1]$ has positions $p_J=[2,3]$.
Note that a position within an axis group differs from its position
in the full tensor: $j_1$ is at position $1$ within $\vec J$, but
position $3$ in $T$.

\paperparagraph{Composing mappings.}
When an operation selects part of a mapped axis group, compilation must
translate positions within the group into positions in the tensor.
Composing mappings performs this translation. A selection $h$ within
$\vec G$ has positions $\gath(p_G,h)$ in tensor $t$, because
\begin{equation}
 \gath(\gath(\operatorname{Ord}_t,p_G),h)
 =\gath(\operatorname{Ord}_t,\gath(p_G,h)).
 \label{eq:gather-composition}
\end{equation}
Both selectors must be valid. In the example above, selecting $j_1$
uses $h=[1]$, giving
$\gath(p_J,h)=\gath([2,3],[1])=[3]$ in $T$.
This is important for TEPL because composed operations may use different
symbolic axis groups. Composing $\gath$ expressions maps positions between
these groups, allowing TEPL to track axis positions in the full tensor.

\subsection{Symbolic Axis Partitioning}
\label{sec:symbolic-axis-partitioning}
When a consumer contracts axes drawn from two producer groups, how do we find
each axis's partner in the other operand? Symbolic axis partitioning answers
this by splitting each group into selected and remaining axes, preserving the
selection order and each selected axis's partner.
Let $c_L,c_R$ be the consumer's contraction-position lists in its left
and right operands, respectively. Continuing the example in
Section~\ref{sec:symbolic-axis-mapping},
Figure~\ref{fig:axis-partitioning} shows a consumer contracting positions $c_L=[3,0]$ in $T$, paired
with $c_R=[2,0]$ in another operand $Z$. The selected axes $[j_1,i_0]$
cross the producer groups, so each group must be split into selected and
remaining axes.

\paperfigure{1}{figures/Symbolic Axis Partitioning (Dot Reassociation).pdf}{Symbolic axis partitioning for unbatched dot reassociation. Splitting selected axes by producer origin preserves their contraction partners. Here $L$ is the selected contraction group, and $F=[f]$ is $Z$'s free group.}{fig:axis-partitioning}{Three panels show the intermediate tensor T with axes i0, i1, j0, j1; selection of positions 3 and 0; partition into selected and remaining I and J subgroups; and preservation of partners at positions 2 and 0 of Z. A table distinguishes positions within producer groups from slots in the contraction list and corresponding partner positions.}

The same selected axis has a different position in each of three lists, so
for each producer group $\vec G\in\{\vec I,\vec J\}$ the compiler records
three position lists:
\begin{itemize}
\item $v_G$: positions of selected axes within $\vec G$;
\item $u_G$: positions of those axes in the consumer's contraction list $c_L$;
\item $z_G$: positions of their paired axes in $Z$.
\end{itemize}
Finding $u_G$ and $v_G$ together is the job of $\mpo$. For lists $A,B$ without
repeated entries, $\mpo(A,B)=(u,v)$, where $u$ lists, in $A$'s order, the
positions in $A$ of the entries that also occur in $B$, and $v$ lists the
positions of those same entries in $B$, so $A_{u_k}=B_{v_k}$. For instance,
$\mpo([3,0],[0,1])=([1],[0])$: the only common entry, $0$, sits at position
$1$ of the first list and position $0$ of the second (full definition in
Appendix~\ref{app:validation-proofs}).
The compiler computes
\begin{equation}
 (u_G,v_G)=\mpo(c_L,p_G),\qquad z_G=\gath(c_R,u_G).
 \label{eq:symbolic-partition-selectors}
\end{equation}
The selected subgroup is $\gath(\vec G,v_G)$; the ordered complement of
$v_G$ selects the remainder. These subgroups are disjoint and cover $\vec G$,
with selected axes following contraction order and remaining axes preserving
their order in $\vec G$. Concatenating the two subgroups therefore need not
restore $\vec G$'s original order.

For the figure ($c_L=[3,0]$, $c_R=[2,0]$, $p_I=[0,1]$, $p_J=[2,3]$):

\begin{center}
\small
\begin{tabular}{@{}lcccc@{}}\toprule
Group & Selected axis & $v_G$ (in $\vec G$) & $u_G$ (in $c_L$) & $z_G=\gath(c_R,u_G)$ (in $Z$)\\\midrule
$\vec I$ & $i_0$ & $[0]$ & $[1]$ & $[0]$\\
$\vec J$ & $j_1$ & $[1]$ & $[0]$ & $[2]$\\\bottomrule
\end{tabular}
\end{center}

The two lists must not be confused: $i_0$ is at position $0$ in $\vec I$ but
position $1$ in $c_L$, and $\gath(c_R,v_I)=\gath(c_R,[0])$ would select
position $2$ in $Z$, the partner of $j_1$, not of $i_0$. Using $u_G$ gives the
correct partners, so the replacement attributes preserve the original
contraction pairings.

\subsection{Symbolic Attribute-Program Construction}
\label{sec:symbolic-program-construction}
Using the mappings and partitions from
Sections~\ref{sec:symbolic-axis-mapping}--\ref{sec:symbolic-axis-partitioning},
the compiler translates Section~\ref{sec:mathematical-synthesis}'s three
passes into one metadata program. Its expressions determine retained axes,
check ordering constraints, and compute RHS axis attributes. Shared
axis-group computations are represented once and reused across attributes.
This program forms the compiler's metadata intermediate representation (IR).

Compilation generates these expressions without requiring concrete ranks
or axis lists. For example, the selected subgroup $\gath(\vec J,v_J)$ has
positions $\gath(p_J,v_J)$ in $T$. At a match, $p_J=[2,3]$ and $v_J=[1]$
evaluate this expression to $[3]$. Each match therefore evaluates the
generated program rather than repeating symbolic inference. The program
operates on match metadata, not tensor values.

The following theorem establishes that the generated feasibility checks
and attribute computations implement Section~\ref{sec:mathematical-synthesis}'s
$\operatorname{check}$ and $\operatorname{build}$ functions within the stated
fragment.

\begin{theorem}[Symbolic compilation correspondence]
\label{thm:symbolic-compilation}
For every admitted semantic plan in Section~\ref{sec:formal-scope}, compilation
constructs a finite metadata program $(E_q,E_f)$ independently of concrete
ranks and the written guard. For every legal metadata instance $m$ in that
section's completeness domain,
\begin{align}
 \lbrack\!\lbrack E_q\rbrack\!\rbrack_m&=\operatorname{check}(m),
 \label{eq:emitted-check}\\
 \operatorname{check}(m)&\Rightarrow
 \lbrack\!\lbrack E_f\rbrack\!\rbrack_m=\operatorname{build}_{\tau}(m).
 \label{eq:emitted-build}
\end{align}
Here $\tau$ fixes unconstrained paired orders using the left operand's
order. The check returns boolean for every legal match in the stated domain.
If it succeeds, the attribute program computes the RHS attributes.
\end{theorem}
\begin{proof}[Proof sketch]
The compiler translates each step of Section~\ref{sec:mathematical-synthesis}
into metadata operations that compute the same axis sets, ordering checks,
and positional attributes. Appendix~\ref{app:symbolic-compilation-proof}
proves this correspondence and shows that program construction terminates
without concrete metadata, the written guard, or successful simplification.
\end{proof}

Together with Theorem~\ref{thm:complete-synthesis}, this establishes
completeness of the generated feasibility/builder program on the stated
domain. Conditional validation checks its use in an executable rewrite.
Appendix~\ref{app:batched-reassociation-program} gives a complete attribute
program for reassociation under explicit sufficient conditions.

\subsection{Implicit Applicability and Acceptance}
\label{sec:compilation-scope}
TEPL compiles the written condition $C$, generated feasibility check $E_q$,
and RHS attribute program $E_f$ once per rule. A legal match first evaluates
$C$, then $E_q$. If either is false, the rule does not apply. The program
$E_f$ runs only after both succeed. Runtime evaluates metadata programs;
it does not repeat synthesis or symbolic proof search.

\paperparagraph{Conditional validation.}
The compiler establishes that $E_q$ reads only source metadata and is a
defined Boolean whenever evaluation reaches it. It then validates $E_f$
under source legality, $C$, and a true result from $E_q$. The obligations
are attribute definedness, replacement legality, and equality of values
and ordered shape. Other construction obligations remain goals to prove.
The compiler does not require $C$ to imply $E_q$. Unresolved validation
returns \code{Unknown}; a false runtime check is ordinary non-application.

\paperparagraph{Acceptance guarantee.}
Under the operation and execution assumptions in
Section~\ref{sec:implementation-scope}, an accepted rule produces a legal,
equivalent replacement on every legal match passing both conditions.
Section~\ref{sec:candidate-validation} and
Appendix~\ref{app:validation-proofs} give the validation argument.
Soundness allows nonnegative extents. Only within the completeness domain
does a false $E_q$ establish that no valid RHS axis assignment exists.

A \emph{sound schema} is a rewrite pattern with a separately proved
attribute construction (Appendix~\ref{app:axis-extensions}). Its conditions
and agreement of supplied fields with constructed fields must hold under
the effective applicability condition. Outside the complete constructor's
fragment, a generated check need only be sufficient; it is not claimed to
accept every feasible match. Table~\ref{tab:guarantee-map} summarizes the
guarantees and their domains.

\begin{table}[!htbp]
\caption{Guarantees and their scope.}
\label{tab:guarantee-map}
\small
\begin{tabularx}{\linewidth}{@{}lY@{}}\toprule
Result & Guarantee\\\midrule
Constructor (Theorem~\ref{thm:complete-synthesis}) & Complete under D1--D4, read-once components, exact arithmetic, and extents $\ge2$; all RHS axis fields unknown.\\
Compilation (Theorem~\ref{thm:symbolic-compilation}) & Shared metadata IR realizes the constructor for admitted plans: $E_q$ implements \code{check}, and $E_f$ implements \code{build} when feasible.\\
Runtime application (Section~\ref{sec:compilation-scope}) & Evaluate $C$ and $E_q$ before $E_f$; false conditions skip the match. Accepted rules preserve values and ordered shape under the execution assumptions.\\
Validation (Theorem~\ref{thm:validator-soundness}) & Establishes safe evaluation of $E_q$ and sound completed candidates under $C$ and $E_q$, including supplied fields; extents $\ge0$.\\
Sound schemas (Lemma~\ref{lem:certified-extensions}) & Sound broadcast/reverse, product, negation, and unit-view schemas under their conditions, including checked supplied fields.\\
Other extensions & Broader reasoning outside the proved fragments has empirical support only.\\\bottomrule
\end{tabularx}
\end{table}
